\documentclass{jsarticle}
\usepackage[dvipdfmx,final]{graphicx}
\usepackage{amsmath}
\usepackage{amssymb}
\begin{document}
\title{素数分布論と素粒子、ゼータ関数　\\
超対称性理論}
\author{Masaaki Yamaguchi}
\date{}
\maketitle

   \includegraphics[width=7cm,clip]{text.png}
   $$ (u + d) + c = e^{-f} - e^{f}, b + (w + s) = e^{-f} + e^{f} $$
   $$ \int \int{1 \over x^s}\mathrm{dvol} - \int \log x \mathrm{dvol}
   = {d \over df}\int\int{1 \over (y \log y)^{1 \over 2}}dy_m
   - \int C dx_m $$
   $$ = e^{-f} - e^{f} $$
   $$ {d \over df}F = {d \over df}\int \int{1 \over (x \log x)^2}dx_m 
   + {d \over df}\int\int{1 \over (y \log y)^{1 \over 2}}dy_m $$
   $$ = e^{f} + e^{-f} $$

     素粒子方程式は、２種類の素粒子と１種類の素粒子の組み合わせで、
   Jones多項式をペアーで構成されている。Rico level理論をベースに
   して、強い力と弱い力、電磁気力が、誤差をD-brane間を重力が
   gravitonとして伝わっているのが、洩れて、反重力をこの
   Weak electric theoryに加えると、全部の力が合わさって、
   統一場理論が完成される。
   これが、ヒッグス場とオイラーの定数を多様体を次元として、
   加群すると、ゼータ関数になる。これが、微分幾何の量子化とガンマ関数
   でもある。



   \includegraphics[width=6cm,clip]{assemble.png}


   $$ ||ds^2|| = e^{-2 \pi T|\psi|}[\eta + \bar{h}(x)]dx^{\mu}dx^{\nu}
   + T^2 d^2 \psi $$
   
   \includegraphics[width=6cm,clip]{energy.png}


      素数分布論が、ゼータ関数を開集合として同型でもあり、ノルム空間、経路積分
   として、位相差が原子間力にもなっている。このエネルギー分布が、
   素数と素粒子方程式、宇宙と異次元のフェルミオンとして、一致している。
   このエネルギー分布が、複素多様体である。
   $$\square = 2 (\cos (i x \log x) + i \sin(i x \log x) $$ 


      宇宙と異次元が、双対性として、超対称性の素粒子を形成している。　

      [reference Lisa Randall, Toshihide Masuoka, Makoto Kobayashi]

   \end{document}
